\subsection{Receptor interventions}

Each receptor is a rigid DOCKSTRING PDBQT file with its search box, a cube of 30~\AA{} centered on the crystallographic ligand. We classified standard residues by the smallest distance $d$ between the box center and a side-chain heavy atom. Contact residues have $d \le 8$~\AA{} and at least two side-chain heavy atoms beyond C$\beta$, shell residues lie at $10 \le d \le 14$~\AA{}, and far residues have all atoms outside the search box enlarged by 8~\AA{}. An intervention removes every atom of a residue except the backbone atoms N, CA, C, O, OXT and the backbone hydrogens and CB, which converts the side chain to alanine. All other lines of the file stay byte-identical, and the graph of the edited receptor carries the residue as ALA. Glycine, alanine, proline, cofactors and metal ions were not edited, and neither were serine and cysteine, which have fewer than two side-chain atoms beyond CB. For each target the plan holds up to 12 single contact truncations chosen evenly over the distance-ranked list, 3 shell truncations and 1 far truncation drawn at random, which serve as negative controls, and 8 multi-residue truncations that combine 2 to 4 of the contact residues that also have single truncations. The seed of each target is $20261001$ plus the sum of the character codes of its name. The receptors were generated from the published files with fixed seeds, and their manifests were hashed so that every pod regenerated identical variants.

\subsection{Re-docking protocol and validation}

Ligand preparation and docking settings are given in the subsection Ligand preparation and docking settings. Every receptor variant of a target was docked with the same ligand files, which gives exactly matched pairs. The measured change of a ligand is $\dS = S_\mathrm{mut} - \overline{S}_\mathrm{wt}$ in \kcal{}. Ligands whose wild-type replicates differed by more than 1.0 \kcal{} in standard deviation, or whose wild-type score was positive, were removed from the truth table, and changes were clipped to $\pm 4$ \kcal{}.

A pilot on ABL1, MAOB and PTGS2 compared engines on identical ligand files (96 ligands docked with Uni-Dock, 24 per target compared with AutoDock Vina 1.1.2). Uni-Dock agreed with AutoDock Vina 1.1.2 (Pearson $r$ = \pilotRLo{} to \pilotRHi{}, mean absolute error \pilotMaeLo{} to \pilotMaeHi{} \kcal{}, bias below 0.1 \kcal{}). At full scale, our preparation reproduced the published DOCKSTRING wild-type scores of the same ligands for the \nParity{} training and development targets (Table~\ref{tab:si_targets}, last column) with a median over targets of \parityR{} for the Pearson $r$ (range \parityRMin{} for \parityRMinTarget{} to \parityRMax{}, \nParityBelow{} targets below 0.9), of \parityMae{} \kcal{} for the mean absolute error and of \parityBias{} \kcal{} for the bias. The nine test targets were not compared, which keeps their labels sealed. An independent check of the measured changes used AutoDock Vina 1.1.2 on the two single and two multi-residue truncations with the largest measured effect on each of three training targets (ESR2, NR3C1 and THRB; \vinaEdits{} edits, \vinaLigPerEdit{} ligands each, \vinaN{} ligand-edit pairs after filtering). The edits were chosen for large effects and the edit means are pooled over three targets, so that the edit-mean correlation of \vinaEditR{} reflects the spread between edits, and the agreement of single ligands is \vinaR{} over all pairs and \vinaWithinR{} within edits.

\subsection{DockMut-1}

The plan contains \nPlanned{} dockings: 73,920 for the 19 held-out targets (10 development and 9 test targets, 160 shared development ligands per target) and 231,424 for the 38 training targets (256 ligands per target drawn from the training pool). Ligands were selected by heavy-atom count in quartiles from 15 to 40 atoms, using SMILES only. Of the \nDock{} scores, \nScoresWt{} belong to wild-type replicates and \nScoresChange{} to edited receptors, \nChangesKept{} measured changes remain after the filters above, and \nNoScore{} planned dockings returned no score. The primary outcome is the set of \nPrimaryPairs{} contact ligand-edit pairs of the \nPrimaryEdits{} contact edits on the 19 held-out targets. Table~\ref{tab:data} gives the composition of the scores obtained.

\begin{table}[t]
\caption{Composition of DockMut-1. Edited receptors are those with at least 20 valid ligands. Dockings are the scores obtained, wild-type replicates included.}
\label{tab:data}
\begin{tabular}{lrrrrrr}
\toprule
Targets & n & Single contact & Multi contact & Controls & Ligands & Dockings \\
\midrule
Training & 38 & 372 & 304 & 152 & 256 & 231{,}319 \\
Development & 10 & 88 & 80 & 40 & 160 & 38{,}068 \\
Test & 9 & 91 & 72 & 34 & 160 & 35{,}840 \\
All & 57 & 551 & 456 & 226 &  & 305{,}227 \\
\bottomrule
\end{tabular}
\end{table}


\subsection{Interventional training}

Each update draws one training target, a batch of 256 wild-type ligands with their DOCKSTRING scores, and six truncated receptors of that target evaluated on a shared set of 96 of its ligands. The shared set exposes every ligand to every sampled edit, which lets the model separate a ligand main effect from an edit-specific effect. The wild-type term is a Huber loss on standardized scores. The intervention term is a Huber loss between the predicted change $S(l,P_\mathrm{mut}) - S(l,P_\mathrm{wt})$ and the measured change in standardized units. Models were trained for 20,000 updates with AdamW (weight decay $10^{-4}$, cosine decay of the learning rate) and the final checkpoint was evaluated without selection. Fifteen percent of the truncated receptors of each training target were withheld to monitor generalization to new edits of known pockets. Multi-residue truncations combine residues that also have single truncations, so a withheld edit can share residues with the training edits of its target, and \shareMultiTrain{}\% of the contact edits are multi-residue truncations. Wild-type ligands are drawn with replacement, the six edits of an update are drawn from all non-withheld edits of the target (including shell and far controls), and the learning rate decays to 5\% of its initial value. The seed fixes the order of the data, the dropout and the label shuffles, and it does not fix the initial weights. The architecture, the weight $\lambda$ of the intervention term and the learning rate were chosen on two validation folds inside the training targets by the registered rule (Table~\ref{tab:si_val} and the registered analysis plan below). The rule selected ResidueSum-Geo with $\lambda = 3$ and a learning rate of $5\times10^{-4}$, and these settings apply to all interventional conditions.

Controls keep the architecture and the wild-type data fixed and change only the intervention labels. The wild-type-only control drops the intervention loss. The ligand-permuted control shuffles the measured changes across the ligands of each receptor, which preserves the mean change of every edit and destroys ligand-specific response. The edit-deranged control gives each edit the measured row of another edit of the same target, which destroys edit-specific information and keeps ligand main effects, and the all-deranged control permutes all measured changes of an update over edits and ligands and keeps only their distribution. A constant-pocket control gives every receptor one constant synthetic pocket, so that its predicted response to an edit is zero by construction. The first implementation of the edit-deranged and all-deranged controls drew uniform random permutations, which return an edit to its own row with probability $1/k$ for $k = 6$ edits per update. We found this after the held-out results were known and trained corrected controls that use a derangement of the edit rows (a random cyclic shift of a random edit order) and, for the all-deranged control, independent permutations of the ligands within each row. Both implementations are reported and marked v1 and v2. All conditions were trained with three seeds, and five conditions (wild type only, measured labels, the ligand-permuted control and the two corrected controls) received three further seeds after the registered outcomes were known, which gives six.

\subsection{Metrics of the response to edits}

For each held-out target and receptor the predicted change is the model score of the edited pocket minus that of the wild-type pocket, clipped to the range of the training labels ($\pm 4$ \kcal{}). Faithfulness is reported as the Pearson correlation over all ligand-receptor pairs of contact truncations, as the correlation of the receptor-mean change across receptors, and as the mean within-receptor correlation after removing the receptor mean, which isolates ligand-specific responses. Table~\ref{tab:main} pools receptors over all targets, whereas the levels of Figure~\ref{fig:levels} and Table~\ref{tab:si_levels} are computed within each target and averaged. The regression gain of the measured on the predicted change tests calibration. Detection of edited residues is the area under the curve for separating contact truncations from shell and far controls by the mean absolute predicted change. The noise ceiling follows from the wild-type replicates, since a measured change has reliability $1 - \sigma_\Delta^2 / \mathrm{Var}(\dS)$ with $\sigma_\Delta = \sigma \sqrt{1 + 1/n}$ for replicate standard deviation $\sigma$ and $n$ wild-type seeds. Intervals come from bootstrap resampling of targets, since targets are the unit of generalization, and differences between conditions use the same resampled targets.

Pooled correlations mix differences between targets, between the edits of a target and between ligands. We therefore also report three levels separately for each held-out target and average them over targets: the correlation over all contact pairs of the target, the correlation of the receptor means of its contact edits, and the mean within-receptor correlation. The shares are estimated per target with the noise of each component subtracted (mutant noise for the edit means and the residuals, and the noise of the shared wild-type mean for the ligand component), negative estimates are set to zero, and the shares are averaged over targets, so that they sum to \shareSum{}\%. The ceiling of the last level is $\sqrt{(v_L + v_I)/(v_L + v_I + v_N)}$ for the shares of the variance of measured changes that belong to the ligand main effect, the edit-by-ligand interaction and replicate noise. To locate the origin of the response, the residue terms of the geometry networks were removed at inference, which keeps the ligand-only term and the global descriptor term. The hot-spot ranking of a model is the Spearman correlation, over the single contact truncations of a target, between the mean absolute predicted change and the mean absolute measured change.

\subsection{Virtual alanine scan}

For the virtual alanine scan, every standard residue with a side chain whose nearest side-chain heavy atom lies within 14~\AA{} of the box center was truncated in silico (two of the residues are serines with a single side-chain heavy atom, which no training edit covers), the graph of each edited pocket was built as for the training edits, and every ligand was queried against every edited pocket. The target of the demonstration was fixed before the scan as the held-out target with the median per-target correlation of the ensemble of the three registered seeds.

\subsection{Descriptor baselines and reference surrogates}

The baseline surrogates are dual encoders (848,002 parameters) trained in the preceding analysis on the same split with 40,000 training ligands per target (39 training targets, against 38 for the networks of this study), a binary cross-entropy loss on top-1\% hits plus an auxiliary Huber loss on scores, and three seeds each for four training conditions: the plain loss, a within-target margin ranking loss, a cross-target reversal ranking loss, and a constant synthetic pocket. The regression head provides the predicted score for the edits, and the hit-probability head, whose minus logit ranked ligands in the preceding analysis, is evaluated as a second readout. Descriptor baselines use no pocket graph. A linear size model uses the heavy-atom count of the ligand, the number of removed atoms and their product. A ridge model and a gradient-boosting model use eight ligand descriptors computed with RDKit and ten edit descriptors, among them the number of removed atoms, the number of edited residues, the distance of the edited residues to the box center, the removed atoms by residue class, and five product features that combine the ligand size with the removed atoms, their distance to the box center and the enclosed volume of the pocket. Two further variants add the ten pocket descriptors of the wild-type receptor and their change under the edit. All baselines were fitted on the training targets that the neural models use. A post hoc variant adds two quantities from the wild-type docking pose of each ligand, the number of ligand heavy atoms within 4.5~\AA{} of the removed side-chain atoms and their smallest distance. This reference needs one docking per ligand, so it is not a surrogate in the sense of this study. It shows how much of the change depends on the pose.
